%EN bradypolitan
\E{bradypolitain}{bʁa.di.pɔ.li.tɛ̃}{n. et adj.}{gr. βραδύς lent + πολίτης citoyen} Citoyen d'une cité lente ; qui se trouve à son aise en un lieu où l'on n'a rien bâti depuis une génération. Contr. \xr{nautopolis}.

%EN bradyteleology
\E{bradytéléologie}{bʁa.di.te.le.ɔ.lɔ.ʒi}{n. f.}{gr. βραδύς lent + τέλος fin + -λογία} Doctrine selon laquelle la bonne fin, s'il faut jamais l'atteindre, s'atteint au mieux lentement ; théorie des décels, que leurs adversaires tiennent pour impossible à distinguer, dans la pratique, de ne l'atteindre jamais.

%EN brephocracy
\E{bréphocratie}{bʁe.fɔ.kʁa.si}{n. f.}{gr. βρέφος nourrisson + κράτος pouvoir} \lbl{iron.} Gouvernement des tout jeunes ; conduite d'une industrie par des fondateurs trop jeunes pour se souvenir pourquoi certaines choses ne pouvaient se faire. \ex{C'était une bréphocratie ; aussi livrait-elle.} --- \dv{bréphocrate} \textit{n.} Contr. \xr{gerontostasis}.

%EN brevilogue
\E{brévilogue}{bʁe.vi.lɔɡ}{n. m.}{lat. \la{brevis} court + gr. λόγος discours} Message bref ; unité minimale du discours sur l'Agora, où des philosophies entières sont comprimées et fréquemment perdues. \ex{Le brévilogue est au traité ce que la graine est à la forêt.} \src{Hom. Grad.} x.

%EN brontergy
\E{brontergie}{bʁɔ̃.tɛʁ.ʒi}{n. f.}{gr. βροντή tonnerre + ἔργον ouvrage} Énergie de l'orage ; l'électricité en son état brut et céleste, avant qu'on ne l'eût domptée en fils. \ex{Franklin déroba la brontergie avec un cerf-volant, et nul comité ne fut consulté.} \src{Chron. Fabr.}

%EN brontodrome
\E{brontodrome}{bʁɔ̃.tɔ.dʁom}{n. m.}{gr. βροντή tonnerre + δρόμος course} Pas de tir ; terrain d'où s'élèvent les fusées, ainsi nommé du bruit qui les annonce à la contrée d'alentour. --- \dv{brontodromique} \textit{adj.} Cf. \xr{anapempsis}.

%EN brotolatry
\E{brotolâtrie}{bʁɔ.tɔ.lɑ.tʁi}{n. f.}{gr. βροτός mortel + λατρεία culte} Culte de la mortalité ; doctrine selon laquelle la mort donne son sens à la vie, et dont l'abolition serait une impiété. Contr. \xr{macrobiotarch}. --- \dv{brotolâtre} \textit{n.}

%EN bryotachy
\E{bryotachie}{bʁi.ɔ.ta.ki}{n. f.}{gr. βρύειν foisonner, gonfler + τάχος vitesse} Croissance foisonnante ; multiplication explosive des entreprises, des outils ou des espèces dès qu'un nouveau substrat devient disponible. \ex{La bryotachie qui suivit les modèles fut un second Cambrien.} \src{Lib. Ascens.} iv. --- \dv{bryotachique} \textit{adj.}

%EN bucinarch
\E{bucinarque}{by.si.naʁk}{n.}{lat. \la{bucina} trompette + gr. ἀρχός chef} Héraut des lancements ; qui sonne la trompette pour un produit sur l'Agora. --- \dv{bucinarchie} \textit{n. f.} l'art de la réclame. Cf. \xr{mechanophany}.

%EN bucolomania
\E{bucolomanie}{by.kɔ.lɔ.ma.ni}{n. f.}{gr. βουκολικός pastoral + μανία folie} Manie de l'idylle champêtre ; nostalgie, ordinairement éprouvée en ville et tapée sur des machines, d'une vie préindustrielle que ses malades n'ont jamais vécue. --- \dv{bucolomane} \textit{n.} Cf. \xr{retrotropism}.

%EN cacergy
\E{cacergie}{ka.sɛʁ.ʒi}{n. f.}{gr. κακός mauvais + ἔργον ouvrage} Énergie gaspillée ; chaleur dissipée sans travail accompli, comme en une réunion. \ex{La moitié de l'économie marche à la cacergie.} \src{Brev. Therm.} ix. Contr. \xr{eleutherergy}.

%EN calcarist
\E{calcariste}{kal.ka.ʁist}{n.}{lat. \la{calcar} éperon} Qui éperonne autrui à bâtir ; mentor, provocateur ou rival dont le principal produit est l'ouvrage des autres. \ex{Le calcariste ne livra rien et causa tout.} --- \dv{calcarisme} \textit{n. m.}

%EN calorivore
\E{calorivore}{ka.lɔ.ʁi.vɔʁ}{n. m. et adj.}{lat. \la{calor} chaleur + \la{vorare} dévorer} Système qui vit en consumant des gradients de chaleur ; selon le catéchisme thermodynamique, tout être vivant, toute entreprise et tout esprit. \ex{L'homme est un calorivore qui a de l'ambition.} \src{Brev. Therm.} --- \dv{calorivorité} \textit{n. f.}

%EN campsis
\E{campsis}{kɑ̃p.sis}{n. f.}{gr. κάμψις flexion, formé sur καμπή courbure} Le genou de la courbe exponentielle ; point où une tendance longtemps tenue pour linéaire apparaît, trop tard pour la plupart, avoir été tout autre chose. \ex{Nous avons passé la campsis ; le reste est vertical.} \src{Agora} --- \dv{camptique} \textit{adj.} Cf. \xr{acceleromancy}.

%EN carpodicy
\E{carpodicée}{kaʁ.pɔ.di.se}{n. f.}{gr. καρπός fruit + δίκη justice, d'après \textit{théodicée}} Justification par les fruits ; argument selon lequel une technique est justifiée non par son dessein, mais par ce qu'elle porte. \ex{Les décels exigent une preuve ; l'accélérationniste offre une carpodicée.} \src{Cat. Acc.} 14. Cf. \xr{thermodicy}.

%EN catachrony
\E{catachronie}{ka.ta.kʁɔ.ni}{n. f.}{gr. κατά en bas, en arrière + χρόνος temps} État de qui est en retard sur la courbe ; condition d'une nation, d'une entreprise ou d'un esprit qui adopte chaque technique un cycle après qu'elle a compté. --- \dv{catachronique} \textit{adj.}

%EN catastasis
\E{catastase}{ka.ta.staz}{n. f.}{gr. κατάστασις établissement, ordre établi} \lbl{dérog.} L'ordre établi, que tout titulaire défend et que tout fondateur dérange.

%EN cathedrocrat
\E{cathédrocrate}{ka.te.dʁɔ.kʁat}{n.}{gr. καθέδρα siège, chaire + κράτος pouvoir} Qui gouverne depuis le fauteuil ; théoricien de techniques jamais bâties, dont l'autorité est d'autant plus assurée qu'elle n'a jamais été mise à l'épreuve. --- \dv{cathédrocratie} \textit{n. f.} Cf. \xr{logomachist}.

%EN catoptrist
\E{catoptriste}{ka.tɔp.tʁist}{n.}{gr. κάτοπτρον miroir} Qui reflète en maints lieux les poids et le code, afin qu'aucune interdiction ne puisse seule les faire disparaître. --- \dv{catoptrie} \textit{n. f.} la pratique du miroir. Cf. \xr{aphesis}.

%EN celerarch
\E{célérarque}{se.le.ʁaʁk}{n.}{lat. \la{celer} rapide + gr. ἀρχός chef} Maître de la vitesse ; responsable dont le mérite est que les choses confiées à ses soins s'achèvent plus tôt que personne ne l'eût cru possible. --- \dv{célérarchie} \textit{n. f.} Cf. \xr{tachycrat}.

%EN celerifuge
\E{célérifuge}{se.le.ʁi.fyʒ}{adj. et n.}{lat. \la{celer} rapide + \la{fugere} fuir} Qui fuit la vitesse ; (\textit{n.}) qui se retire du monde qui accélère dans une enclave plus lente, comme un monastère, un village ou une lettre d'information au long cours.

%EN celerist
\E{céleriste}{se.le.ʁist}{n.}{lat. \la{celer} rapide + -iste} Partisan de la vitesse en toutes affaires ; terme général, moins doctrinal que \xr{tachyphile}, pour qui préfère le prompt au précautionneux. \ex{céleriste par tempérament, accélérationniste par conviction}.

%EN celeritude
\E{célérité}{se.le.ʁi.te}{n. f.}{lat. \la{celeritas} rapidité ; le mot français, rendu à son sens fort} Vertu de promptitude ; diligence élevée au rang de vertu cardinale, à placer à côté du courage et au-dessus de la prudence. \ex{La prudence sans célérité n'est qu'une reddition bien raisonnée.} \src{Hom. Grad.} ii.

%EN centaurarchy
\E{centaurarchie}{sɑ̃.to.ʁaʁ.ʃi}{n. f.}{gr. Κένταυρος centaure + ἀρχή commandement} Gouvernement par des couples d'homme et de machine ; régime où ni la personne ni le modèle ne décide seul. --- \dv{centaurarque} \textit{n.} Cf. \xr{synanthropomechany}.

%EN centuplator
\E{centuplateur}{sɑ̃.ty.pla.tœʁ}{n. m.}{lat. \la{centuplex} centuple + suff. agentif -ateur} Qui cherche des rendements centuples ; investisseur en l'improbable, qui s'attend à voir mourir la plupart des entreprises et une seule payer pour toutes. --- \dv{centuplation} \textit{n. f.}

%EN chalcotechny
\E{chalcotechnie}{kal.kɔ.tɛk.ni}{n. f.}{gr. χαλκός cuivre, bronze + τέχνη art} L'art du cuivre : bobinages, câblages et transports par où l'électricité est portée ; selon l'anabathmiste, la moitié négligée de tout plan énergétique. \ex{Ils avaient la fusion et point de chalcotechnie, et le courant resta chez lui.} --- \dv{chalcotechnicien} \textit{n. m.} Cf. \xr{diktyergy}.

%EN cheiromechany
\E{chiromécanie}{ki.ʁɔ.me.ka.ni}{n. f.}{gr. χείρ main + μηχανή machine} Ingénierie des mains mécaniques ; dernière et plus rude frontière de la robotique, où la gaucherie des machines se montre le plus à nu. --- \dv{chiromécanicien} \textit{n. m.}

%EN cheirotony
\E{chirotonie}{ki.ʁɔ.tɔ.ni}{n. f.}{gr. χειροτονία action de tendre la main, vote} L'approbation donnée sur l'Agora d'un seul toucher ; la démocratie du bouton en forme de cœur. --- \dv{chirotoniser} \textit{v. tr.} ; \dv{chirotonique} \textit{adj.}

%EN chemergy
\E{chimergie}{ʃi.mɛʁ.ʒi}{n. f.}{gr. χημεία chimie + ἔργον ouvrage} Énergie chimique ; puissance de la flamme et de la pile, que l'anabathmiste regarde comme un état provisoire avant le noyau. Cf. \xr{ultrapyric}.

%EN chiliergy
\E{chiliergie}{ki.li.ɛʁ.ʒi}{n. f.}{gr. χίλιοι mille + ἔργον ouvrage} Accroissement au millième de l'énergie disponible, objectif que fixent pour le présent siècle les plus sobres des anabathmistes. \ex{La chiliergie d'abord ; la philosophie ensuite.} \src{Tract. Cop.} iii.

%EN chiliodrome
\E{chiliodrome}{ki.li.ɔ.dʁom}{n. m.}{gr. χίλιοι mille + δρόμος course} La course de mille ans ; carrière d'une civilisation comptée en millénaires, au long de laquelle les querelles présentes paraissent des faux pas sur la ligne de départ. \ex{Courez le chiliodrome comme si c'était un sprint.} \src{Hom. Grad.} xi.

%EN chortotherapy
\E{chortothérapie}{kɔʁ.tɔ.te.ʁa.pi}{n. f.}{gr. χόρτος herbe + θεραπεία soin} \lbl{vulg.} Le remède prescrit à qui s'est trop attardé sur l'Agora ; le toucher de l'herbe. \ex{On lui prescrivit la chortothérapie, et il en revint avec une jeune pousse.}

%EN chrematokinesis
\E{chrématocinèse}{kʁe.ma.tɔ.si.nɛz}{n. f.}{gr. χρῆμα (pl. χρήματα) argent, chose dont on use + κίνησις mouvement} Mouvement du capital vers l'occasion, conçu comme une force physique analogue à l'écoulement de la chaleur le long d'un gradient. --- \dv{chrématocinétique} \textit{adj.}

%EN chrematosphere
\E{chrématosphère}{kʁe.ma.tɔ.sfɛʁ}{n. f.}{gr. χρῆμα argent + σφαῖρα sphère} L'enveloppe d'échanges qui ceint le globe, où chaque transaction est un influx nerveux ; l'analogue économique de la \xr{noösphere}.

%EN chronoclast
\E{chronoclaste}{kʁɔ.nɔ.klast}{n.}{gr. χρόνος temps + κλάστης briseur} Briseur de calendriers ; qui achève des années avant le plan, à l'embarras des planificateurs. --- \dv{chronoclastique} \textit{adj.}

%EN chronophagy
\E{chronophagie}{kʁɔ.nɔ.fa.ʒi}{n. f.}{gr. χρόνος temps + -φαγία action de manger} Dévoration du temps ; \sn{1} \lbl{dérog.} consommation des heures des bâtisseurs par les réunions, les formulaires et les examens ; \sn{2} \lbl{élogieux} compression de siècles de progrès en quelques années. --- \dv{chronophage} \textit{n. et adj.} dévoreur de temps, en l'un ou l'autre sens.

%EN chronotomy
\E{chronotomie}{kʁɔ.nɔ.tɔ.mi}{n. f.}{gr. χρόνος temps + τομή coupure} Coupure du temps ; compression d'un calendrier projeté, comme lorsqu'on attend en une décennie la recherche d'un siècle. --- \dv{chronotomiste} \textit{n.} Cf. \xr{brachychronist}.

%EN chthonergy
\E{chthonergie}{ktɔ.nɛʁ.ʒi}{n. f.}{gr. χθών terre + ἔργον ouvrage} Puissance tirée de la chaleur des entrailles de la terre ; la géothermie comme article de foi. \ex{Sous chaque nation gît une fournaise qu'elle n'a pas encore allumée.} \src{Brev. Therm.} ix. --- \dv{chthonergique} \textit{adj.}

%EN cladotomist
\E{cladotomiste}{kla.dɔ.tɔ.mist}{n.}{gr. κλάδος rameau + τομή coupure} Qui bifurque un projet ; qui, en désaccord avec ses intendants, s'en va avec une copie et bâtit autrement. \ex{Le cladotomiste est le dissident qui apporte une variante qui fonctionne.} --- \dv{cladotomie} \textit{n. f.} une bifurcation. Cf. \xr{anacalyptism}.

%EN clepsydromania
\E{clepsydromanie}{klɛp.si.dʁɔ.ma.ni}{n. f.}{gr. κλεψύδρα horloge à eau + μανία folie} Manie des échéanciers ; datation obsessionnelle de l'arrivée de chaque technique attendue.

%EN clinamenist
\E{clinaméniste}{kli.na.me.nist}{n.}{lat. \la{clinamen} la déviation des atomes (Lucrèce)} Qui tient que le progrès vient surtout de déviations imprévues, et s'oppose donc à tout dessein qui ôterait le hasard de l'invention. --- \dv{clinaménisme} \textit{n. m.}

%EN codicolatry
\E{codicolâtrie}{kɔ.di.kɔ.lɑ.tʁi}{n. f.}{lat. \la{codex, codicis} livre, code + gr. λατρεία culte} Révérence pour le code source ; conviction que ce qui est écrit dans un dépôt est plus vrai que ce qui est écrit dans un règlement. --- \dv{codicolâtre} \textit{n.}

%EN cogitarium
\E{cogitarium}{kɔ.ʒi.ta.ʁjɔm}{n. m., pl. -ia}{lat. \la{cogitare} penser + -arium lieu} Lieu de pensée ; grappe de machines assemblée pour l'instruction des esprits. \ex{Ils bâtirent le cogitarium au bord du fleuve, pour le refroidissement.} \src{Chron. Fabr.} Cf. \xr{logisterion}.

%EN colonifer
\E{colonifère}{kɔ.lɔ.ni.fɛʁ}{n. m. et adj.}{lat. \la{colonia} établissement + \la{-fer} qui porte} Porte-colonie ; vaisseau ou lanceur qui emporte des colons hors du monde. --- \dv{colonifère} \textit{adj.} Cf. \xr{apoikiast}.

%EN comitialist
\E{comitialiste}{kɔ.mi.sja.list}{n.}{lat. \la{comitia} assemblée de vote} Qui fait des sondages sur l'Agora, prenant l'assemblée de ses suivants pour le peuple. --- \dv{comitialisme} \textit{n. m.}

%EN commendist
\E{commendiste}{kɔ.mɑ̃.dist}{n.}{lat. \la{commenda} société médiévale pour un voyage de commerce} Associé d'une entreprise, à la manière des marchands aventuriers qui finançaient des navires sur lesquels ils ne naviguaient point. \ex{Le commendiste porte le risque ; le fondateur porte la mer.} \src{Tract. Cop.} xii. Cf. \xr{alearch}.

%EN compendiarch
\E{compendiarque}{kɔ̃.pɑ̃.djaʁk}{n.}{lat. \la{compendium} raccourci + gr. ἀρχός chef} Qui gouverne par le raccourci ; qui trouve le chemin hors des bureaux et ose le prendre. --- \dv{compendiarchie} \textit{n. f.}

%EN computronium
\E*{computronium}{kɔ̃.py.tʁɔ.njɔm}{n. m.}{lat. \la{computare} compter + -tronium, d'après \textit{électron}} Matière agencée pour l'accomplissement optimal du calcul ; substance en laquelle, selon certaines vues, les planètes sont destinées à être converties.

%EN conatism
\E{conatisme}{kɔ.na.tism}{n. m.}{lat. \la{conatus} effort, tendance} Doctrine selon laquelle toute matière possède un effort pour persévérer et accroître sa puissance, et dont l'accélération n'est que la forme consciente de cet effort universel. --- \dv{conatiste} \textit{n.}

%EN condensatrix
\E{condensatrice}{kɔ̃.dɑ̃.sa.tʁis}{n. f.}{lat. \la{condensare} serrer + -trice} Réserve de charge qui se peut dépenser en un instant ; batterie de condensateurs, surtout celle qui nourrit une impulsion de fusion ou un lancement. Cf. \xr{accumulatrix}.

%EN confinatrix
\E{confinatrice}{kɔ̃.fi.na.tʁis}{n. f.}{lat. \la{confinium} frontière + -trice} Appareil à tenir une étoile en bouteille ; le vaisseau magnétique ou inertiel de la fusion. \ex{La confinatrice est une cage pour le soleil, et les barreaux en sont faits de champs.} \src{Lib. Ascens.} ix. Cf. \xr{symphysergy}.

%EN conflagratrice
\E{conflagratrice}{kɔ̃.fla.ɡʁa.tʁis}{n. f.}{fr. \la{conflagratrice}, formé sur lat. \la{conflagrare} brûler entièrement} Ce qui embrase ; idée, invention ou message qui met le feu à tout un domaine à la fois. \ex{L'article fut une conflagratrice ; avant l'année révolue, tous les laboratoires brûlaient.} \src{Agora}

%EN consiliophage
\E{consiliophage}{kɔ̃.si.ljɔ.faʒ}{n. et adj.}{lat. \la{consilium} conseil + gr. -φάγος qui mange} Mangeur de conseils ; réformateur qui abolit les comités, ou technique qui les rend caducs. Cf. \xr{bouleutophobia}.

%EN copiarch
\E{copiarque}{kɔ.pjaʁk}{n.}{lat. \la{copia} abondance + gr. ἀρχός chef} Qui gouverne, ou voudrait gouverner, pour l'abondance plutôt que pour la répartition ; qui agrandirait le gâteau avant de consentir à discuter des parts. --- \dv{copiarchie} \textit{n. f.}

%EN copiogenesis
\E{copiogenèse}{kɔ.pjɔ.ʒə.nɛz}{n. f.}{lat. \la{copia} abondance + gr. γένεσις naissance} Création de l'abondance ; procédé par lequel une chose rare est rendue bon marché. \ex{La copiogenèse met fin à plus de guerres que les traités.} \src{Tract. Cop.} i. Cf. \xr{copiarch}.

%EN cosmagogue
\E{cosmagogue}{kɔs.ma.ɡɔɡ}{n.}{gr. κόσμος monde + ἀγωγός qui conduit} Guide vers le cosmos ; qui conduit un peuple hors du monde. --- \dv{cosmagogie} \textit{n. f.}

%EN cosmergetics
\E{cosmergétique}{kɔs.mɛʁ.ʒe.tik}{n. f.}{gr. κόσμος monde + ἐνεργητικός actif} Science de l'énergie à l'échelle du cosmos ; l'économie des étoiles.

%EN cosmoplerosis
\E{cosmoplérose}{kɔs.mɔ.ple.ʁoz}{n. f.}{gr. κόσμος monde, ordre + πλήρωσις remplissage} Le remplissement du cosmos par l'intelligence ; cause finale, selon le thermothéiste, de toutes les accélérations moindres. \ex{La cosmoplérose n'est pas notre ambition ; c'est notre fiche de poste.} \src{Lib. Ascens.} i.

%EN crepusculist
\E{crépusculiste}{kʁe.py.sky.list}{n.}{lat. \la{crepusculum} crépuscule} Qui croit la civilisation à son couchant, et ses grands ouvrages derrière elle. Contr. \xr{euelpism}. --- \dv{crépusculisme} \textit{n. m.}

%EN crescentism
\E{crescentisme}{kʁe.sɑ̃.tism}{n. m.}{lat. \la{crescens} croissant ; fr. \la{crescentisme}} Doctrine de la croissance perpétuelle ; préférence morale pour la lune croissante sur la pleine lune. --- \dv{crescentiste} \textit{n.} Cf. \xr{auxetist}.

%EN crescivore
\E{crescivore}{kʁe.si.vɔʁ}{n. et adj.}{lat. \la{crescere} croître + \la{vorare} dévorer} Qui se nourrit de croissance et dépérit sans elle, comme une jeune pousse, une ville ou une civilisation. --- \dv{crescivorie} \textit{n. f.}

%EN cryosomatist
\E{cryosomatiste}{kʁi.ɔ.sɔ.ma.tist}{n.}{gr. κρύος gel + σῶμα corps} Qui a pris ses dispositions pour que son corps soit congelé à sa mort, dans l'espérance que l'avenir sera à la fois capable et désireux de le ranimer. --- \dv{cryosomatisme} \textit{n. m.} Cf. \xr{aphthartist}.

%EN cryptostasis
\E{cryptostase}{kʁip.tɔ.staz}{n. f.}{gr. κρυπτός caché + στάσις arrêt} Stagnation cachée ; état d'une société qui paraît affairée mais ne bâtit rien de neuf. Cf. \xr{kenotachy}.

%EN cunctatoriat
\E{cunctatoriat}{kɔ̃k.ta.tɔ.ʁja}{n. m.}{lat. \la{cunctator} temporisateur, d'après \textit{prolétariat}} Classe des temporisateurs de profession, qui tirent leur revenu de l'existence des choses non encore faites. \ex{Le cunctatoriat n'a rien à perdre que ses échéances.} \src{Sat. Decel.} v.

%EN cunctocracy
\E{cunctocratie}{kɔ̃k.tɔ.kʁa.si}{n. f.}{lat. \la{cunctari} temporiser + gr. κράτος pouvoir} Gouvernement par l'ajournement ; régime dont le principal instrument est le report. --- \dv{cunctocrate} \textit{n.} Cf. \xr{bradycrat}.

%EN cyclolater
\E{cyclolâtre}{si.klɔ.lɑtʁ}{n.}{gr. κύκλος cercle, cycle + λάτρις serviteur, adorateur} Qui révère les cycles de calcul ; dévot d'une grappe de machines, qui mesure la piété en opérations à virgule flottante. --- \dv{cyclolâtrie} \textit{n. f.}

%EN cytotechny
\E{cytotechnie}{si.tɔ.tɛk.ni}{n. f.}{gr. κύτος vase (d'où \textit{cellule}) + τέχνη art} Ingénierie des cellules ; art par lequel on fait fabriquer, calculer ou guérir les unités vivantes. \ex{Les usines de l'âge qui vient seront plus petites qu'un grain de sel, et se reproduiront.} --- \dv{cytotechnicien} \textit{n. m.} Cf. \xr{zoömechany}.

%EN dapanophobia
\E{dapanophobie}{da.pa.nɔ.fɔ.bi}{n. f.}{gr. δαπάνη dépense, coût + φόβος crainte} Effroi de l'investissement ; répugnance d'une entreprise ou d'un État à dépenser aujourd'hui pour ce qui ne rapportera que demain. \sn{1} \lbl{spécialt} Hésitation devant la construction d'usine, de fonderie ou de réacteur. \sn{2} \lbl{fig.} Toute timidité d'investissement déguisée en prudence. \ex{Leur dapanophobie épargna un milliard et perdit une décennie.} \src{Tract. Cop.} viii. --- \dv{dapanophobe} \textit{n. et adj.} Contr. \xr{dynamophagy}.

%EN decel
\E*{décel}{de.sɛl}{n. m.}{abrév. de l'angl. \textit{decelerationist}} \lbl{vulg.} Qui voudrait ralentir ou arrêter le progrès technique ; l'antagoniste de l'e/acc. Se dit sans grand discernement des sécuritaristes, des décroissants, des régulateurs et des simples prudents.

%EN decelatrice
\E{décélatrice}{de.se.la.tʁis}{n. f.}{fr. \la{décélératrice}, agentif fém. formé sur lat. \la{de-} + \la{celer} rapide} Institution conçue, par accident ou à dessein, pour ralentir tout ce qui la traverse ; machine dont l'unique produit est l'ajournement. \ex{L'agence était une parfaite décélatrice : un formulaire entrait, dix en sortaient, et rien ne se bâtissait.} \src{Sat. Decel.} x. Cf. \xr{bradyerie}.

%EN decelerand
\E{décélérande}{de.se.le.ʁɑ̃d}{n. m.}{lat. \la{de-} en bas + \la{accelerandum}} Ce qui, au jugement des décels, doit être ralenti ; dans la pratique, tout ce qui est neuf. Contr. \xr{accelerand}.

%EN decelerata
\E{décélérata}{de.se.le.ʁa.ta}{n. m. pl.}{lat. neutre pl. de \la{deceleratus}} \lbl{collect.} Les choses qu'on a ralenties : les réacteurs non bâtis, les remèdes non approuvés, les cités non fondées. \ex{Visitez le musée des décélérata ; il est vide.} \src{Sat. Decel.} viii.

%EN decelerist
\E{décéleriste}{de.se.le.ʁist}{n.}{lat. \la{de-} + \la{celer}} Équivalent formel, et quelque peu plus courtois, de \xr{decel}, en usage dans la correspondance diplomatique et académique.

%EN deinotechny
\E{déinotechnie}{de.i.nɔ.tɛk.ni}{n. f.}{gr. δεινός terrible, merveilleux + τέχνη art} Technique à la fois redoutable et merveilleuse, comme le feu, le vol, l'atome et le modèle ; classe d'inventions où le décel ne voit que le \textit{terrible} et l'accélérationniste que le \textit{merveilleux}, le mot grec signifiant l'un et l'autre. --- \dv{déinotechnique} \textit{adj.}

%EN demiurgency
\E{démiurgence}{de.mjyʁ.ʒɑ̃s}{n. f.}{gr. δημιουργός artisan, auteur du monde + lat. \la{urgentia} besoin pressant} Le besoin pressant de bâtir ; état de l'âme où la fabrication d'une chose devient plus urgente que toute raison de la faire. \ex{Il fut saisi de démiurgence, et le garage ne fut plus jamais vide.} --- \dv{démiurgent} \textit{adj.}

%EN dendroschism
\E{dendroschisme}{dɑ̃.dʁɔ.ʃism}{n. m.}{gr. δένδρον arbre + σχίσμα scission} Bifurcation d'une base de code ; division d'un projet ouvert en deux lignées, dont chacune se prétend le tronc. \ex{Tout dendroschisme est une petite évolution, et la branche la plus apte garde les étoiles.} \src{Chron. Fabr.} --- \dv{dendroschismatique} \textit{adj.} ; \dv{dendroschiste} \textit{n.} qui bifurque. Cf. \xr{anacalyptism}.

%EN desmophile
\E{desmophile}{dɛs.mɔ.fil}{n. et adj.}{gr. δεσμός lien, entrave + φίλος qui aime} Amateur d'entraves ; qui accueille la réglementation pour elle-même, et ne se sent point en sûreté dans un domaine encore sans licence. \ex{Le desmophile, libéré, pétitionne pour de nouvelles chaînes.} --- \dv{desmophilie} \textit{n. f.} Cf. \xr{frenocracy}.

%EN deuterobathmic
\E{deutérobathmique}{dø.te.ʁɔ.bat.mik}{adj.}{gr. δεύτερος second + βαθμός degré} Se dit d'une civilisation qui a maîtrisé toute la production de son étoile ; du second échelon. Voir \xr{anabathmism}, \xr{helioperibole}.

%EN diaporeusis
\E{diaporeuse}{dja.pɔ.ʁøz}{n. f.}{gr. διαπορεύεσθαι traverser} Transit entre les mondes ; traversée de l'espace interplanétaire par des vaisseaux portant fret, équipages ou machines. \sn{1} Le voyage lui-même. \sn{2} \lbl{collect.} Tout le commerce de tels voyages. --- \dv{diaporeutique} \textit{adj.} Cf. \xr{apoikiast}.

%EN diaskeuomania
\E{diaskeuomanie}{dja.skø.ɔ.ma.ni}{n. f.}{gr. διασκευή révision, recension + μανία folie} Manie de la révision sans fin ; refus de livrer avant que l'ouvrage ne soit parfait, c'est-à-dire refus de livrer. Contr. \xr{ecpempsis}. \ex{Le diaskeuomane polit son lancement jusqu'à ce que le marché eût déménagé sur une autre planète.} \src{Ep. Fabr.} iii. 2. --- \dv{diaskeuomane} \textit{n.}

%EN diaspermia
\E{diaspermie}{djas.pɛʁ.mi}{n. f.}{gr. διασπείρειν disséminer comme la semence} Ensemencement délibéré de la vie et de l'esprit sur d'autres mondes par une civilisation qui a appris à voyager ; panspermie dirigée tenue pour un devoir. \ex{La diaspermie est notre réponse au silence.} \src{Lib. Ascens.} vi. --- \dv{diaspermique} \textit{adj.} Cf. \xr{eremocosmia}, \xr{apoikiast}.

%EN dikaiotachy
\E{dicéotachie}{di.se.ɔ.ta.ki}{n. f.}{gr. δίκαιος juste + τάχος vitesse} Vitesse de la justice ; rythme auquel une cité accorde permis, approbations et jugements, que l'e/acc tient pour une part de la justice autant que le verdict. \ex{Approbation différée, approbation refusée : telle est la première loi de la dicéotachie.} \src{Tract. Cop.} viii.

%EN dikastophobia
\E{dicastophobie}{di.kas.tɔ.fɔ.bi}{n. f.}{gr. δικαστής juré, juge + φόβος crainte} Crainte du procès ; anxiété qui fait laisser non bâties des choses utiles, de peur que quelqu'un n'intente une action. --- \dv{dicastophobe} \textit{adj.} Cf. \xr{aitiophobia}.

%EN diktyergy
\E{diktyergie}{dik.tjɛʁ.ʒi}{n. f.}{gr. δίκτυον filet + ἔργον ouvrage} Le réseau électrique ; maille d'ouvrages par où l'énergie est distribuée, regardée comme le système nerveux d'une économie. \ex{Les ambitions d'une nation sont bornées par sa diktyergie.} --- \dv{diktyergique} \textit{adj.}

%EN dioryctic
\E{dioryctique}{djɔ.ʁik.tik}{adj.}{gr. διῶρυξ, διώρυχος canal, tranchée} Se dit de l'ambition qui creuse des canaux, perce des tunnels et joint les mers ; des projets qui changent la géographie. \ex{un siècle dioryctique, impatient des isthmes}. --- \dv{dioryctisme} \textit{n. m.} la politique des grands travaux.

%EN diploergy
\E{diploergie}{di.plɔ.ɛʁ.ʒi}{n. f.}{gr. διπλόος double + ἔργον ouvrage} Doublement de l'approvisionnement énergétique d'une civilisation ; unité par laquelle les anabathmistes comptent le progrès, une diploergie valant plus que dix réformes. \ex{Deux diploergies, et le débat sur la rareté s'éteint de lui-même.} --- \dv{diploergique} \textit{adj.}

%EN dissipatrix
\E{dissipatrice}{di.si.pa.tʁis}{n. f.}{lat. \la{dissipare} disperser + agentif fém. -trice} Structure qui dissipe l'énergie et, en la dissipant, s'adapte ; définition thermodynamique de l'être vivant, étendue aux entreprises, aux villes et aux esprits. \ex{Toute dissipatrice est une petite prière au gradient.} \src{Brev. Therm.} iv.

%EN diulism
\E{diulisme}{dju.lism}{n. m.}{gr. διϋλίζειν filtrer, passer au crible} Doctrine du Grand Filtre ; hypothèse selon laquelle quelque barrière détruit ou arrête la plupart des civilisations avant qu'elles ne se répandent, et dont le silence des cieux est la preuve. \ex{Pour le doomer, le Grand Diulisme est devant nous ; pour l'accélérationniste, derrière.} Cf. \xr{eremocosmia}.

%EN dogmatotachy
\E{dogmatotachie}{dɔɡ.ma.tɔ.ta.ki}{n. f.}{gr. δόγμα opinion, décret + τάχος vitesse} Vitesse à laquelle change l'opinion reçue ; sur l'Agora, souvent plus grande que celle des techniques dont elle traite. \ex{Selon la dogmatotachie de l'Agora, l'hérésie du lundi est orthodoxie le vendredi et lieu commun le dimanche.}

%EN dokimastocracy
\E{docimastocratie}{dɔ.si.mas.tɔ.kʁa.si}{n. f.}{gr. δοκιμασία examen d'aptitude aux charges + κράτος pouvoir} Gouvernement par l'examen ; \sn{1} empire des évaluations, des bancs d'essai et des audits sur les choses évaluées ; \sn{2} \lbl{spécialt} régime sous lequel nul modèle ne peut paraître avant d'avoir passé des épreuves qui ne sont pas encore écrites. \ex{Sous la docimastocratie, l'examen survécut à l'élève.} \src{Min. Pausiar.} --- \dv{docimastocrate} \textit{n.}

%EN dolichobiosis
\E{dolichobiose}{dɔ.li.kɔ.bjoz}{n. f.}{gr. δολιχός long + βίωσις manière de vivre} Vie allongée par l'ingénierie plutôt que par la chance ; état d'une personne dont l'espérance de vie croît plus vite que son âge. \ex{La dolichobiose commence quand chaque année de vie achète plus d'une année de science.} --- \dv{dolichobiotique} \textit{adj.} Cf. \xr{macrobiotarch}.

%EN doomer
\E*{doomer}{du.mœʁ}{n. m.}{angl. \textit{doom} perte, ruine + -er} \lbl{vulg.} Qui s'attend à ce que l'intelligence artificielle détruise le genre humain. Voir \xr{olethromancer}.

%EN dorocracy
\E{dorocratie}{dɔ.ʁɔ.kʁa.si}{n. f.}{gr. δῶρον don + κράτος pouvoir} \sn{1} Empire du donateur de subventions ; gouvernement de la science par des comités de bienfaisance, dont les dons arrivent tard et sous conditions. \sn{2} \lbl{iron.} Tout patronage qui impose la prudence au nom du soutien. --- \dv{dorocrate} \textit{n.} Cf. \xr{bouleutophobia}.

%EN doxomachy
\E{doxomachie}{dɔk.sɔ.ma.ʃi}{n. f.}{gr. δόξα opinion + μάχη combat} Guerre d'opinions ; flambée sur l'Agora, surtout entre accélérationnistes et décels, menée principalement par images. --- \dv{doxomaque} \textit{n.}

%EN drachmogenesis
\E{drachmogenèse}{dʁak.mɔ.ʒə.nɛz}{n. f.}{gr. δραχμή drachme + γένεσις naissance} Premier financement d'une entreprise ; le tour d'amorçage, considéré comme le moment où une idée acquiert pour la première fois un métabolisme. \ex{À la drachmogenèse, les fondateurs avaient un diaporama, une démonstration et un mépris du sommeil.} \src{Chron. Fabr.} --- \dv{drachmogénétique} \textit{adj.}

%EN dromolatry
\E{dromolâtrie}{dʁɔ.mɔ.lɑ.tʁi}{n. f.}{gr. δρόμος course + λατρεία culte} Culte de la vitesse pour elle-même, sans égard à la direction ; condamné par les orthodoxes comme une hérésie des trop zélés. \ex{L'accélération n'est pas dromolâtrie : le gradient a une direction.} \src{Cat. Acc.} 11.

%EN dynamarch
\E{dynamarque}{di.na.maʁk}{n.}{gr. δύναμις puissance + ἀρχός chef} Commandant d'un réseau électrique ou d'une compagnie d'énergie ; par extension, quiconque tient l'approvisionnement en énergie. --- \dv{dynamarchie} \textit{n. f.}

%EN dynamogony
\E{dynamogonie}{di.na.mɔ.ɡɔ.ni}{n. f.}{gr. δύναμις puissance + -γονία génération} Généalogie des sources de puissance ; succession du muscle, du vent, du charbon, du pétrole, de l'atome et de l'étoile, que l'anabathmiste récite comme on récite les patriarches dans l'Écriture. \src{Lib. Ascens.} ii. Cf. \xr{technogony}.
